% META: {"id": "1NSI-LANG-CO-006", "chapitre": "1NSI-LANGAGE", "type_objet": "corrige", "exercice_ref": "1NSI-LANG-EX-006", "capacites": ["P-LANG-01"], "mode_creation": "ex_nihilo", "sources_inspiration": [], "status": "generated"}
% BEGIN-VERIFY
% def total_panier(prix):
%     total = 0
%     for i in range(len(prix)):
%         if prix[i] > 0:
%             total = total + prix[i]
%     return round(total * 100) / 100
% assert total_panier([3.5, -1, 2.0, 4.25]) == 9.75
% assert total_panier([]) == 0
% assert total_panier([-2, -8]) == 0
% assert 3.5 + 2.0 + 4.25 == 9.75
% def total_avec_while(prix):
%     total, i = 0, 0
%     while i < len(prix):
%         if prix[i] > 0:
%             total = total + prix[i]
%         i = i + 1
%     return round(total * 100) / 100
% assert total_avec_while([3.5, -1, 2.0, 4.25]) == total_panier([3.5, -1, 2.0, 4.25])
% assert round(0.1 * 3 * 100) / 100 == 0.3
% END-VERIFY
\begin{corrige}{1NSI-LANG-EX-006}
\textbf{1.} Les cinq constructions sont toutes présentes :

\begin{center}
\begin{tabular}{ll}
\hline
construction & occurrence \\
\hline
séquence & les trois instructions du corps, exécutées l'une après l'autre \\
affectation & \lstinline{let total = 0;} et \lstinline{total = total + prix[i];} \\
boucle bornée & \lstinline{for (let i = 0; i < prix.length; i = i + 1)} \\
conditionnelle & \lstinline{if (prix[i] > 0)} \\
appel de fonction & \lstinline{Math.round(total * 100)} \\
\hline
\end{tabular}
\end{center}

\textbf{2.} La traduction ne change que l'habillage :
\begin{python}
def total_panier(prix):
    total = 0
    for i in range(len(prix)):
        if prix[i] > 0:
            total = total + prix[i]
    return round(total * 100) / 100
\end{python}

\textbf{3.} Sur \lstinline{[3.5, -1, 2.0, 4.25]}, la fonction additionne
$3{,}5+2{,}0+4{,}25=9{,}75$. Sur la liste vide comme sur
\lstinline{[-2, -8]}, elle renvoie $0$ : la boucle ne fait aucun tour dans le
premier cas, et la conditionnelle est fausse à chaque tour dans le second.

Pour une caisse, ces deux zéros ne veulent pourtant pas dire la même chose : un
panier vide est normal, un panier composé uniquement de montants négatifs est
un remboursement, et l'afficher comme un total de $0$ euro serait une erreur
commerciale. La fonction est correcte au sens du calcul, insuffisante au sens du
métier — c'est une distinction qu'aucun test de valeur de retour ne fera à notre
place.

\textbf{4.} C'est une différence de \textbf{syntaxe}. Les deux codes emploient
exactement les mêmes constructions, dans le même ordre ; seule change la façon
de les écrire — accolades et point-virgules d'un côté, indentation de l'autre.
C'est précisément ce que veut dire « constructions communes aux langages de
programmation » : ce qui se traduit d'un langage à l'autre sans rien
réorganiser.

\textbf{5.} La \textbf{boucle non bornée} (\lstinline{while}) est absente : ici
le nombre de tours est connu d'avance, c'est le nombre d'articles. La caisse en
aurait besoin, par exemple, pour redemander un code-barres tant que la saisie
est invalide : on ne sait pas à l'avance combien de fois le client s'y
reprendra. On peut d'ailleurs réécrire la même fonction avec un
\lstinline{while} et un compteur explicite, et vérifier qu'elle donne le même
total — mais rien n'y gagne en clarté, car le nombre de tours, lui, reste connu.
\end{corrige}
